\section{Concentration of the original source at criticality and zero residence}
\label{sec:v87-original}
The finite-count envelope can be inserted directly into the scalar
source-shift argument. The exceptional contribution can then be kept as
its original weighted starting mass, instead of enlarged to an optical
potential. No mass is removed from the original event.

Let $\ell(x)=\tau^-(x)+\tau^+(x)$ denote the total roof residence of
its maximal regular exact-event segment. The Borel parametrization gives
$\ell=b_j-a_j$. For $a(h),\delta(h)$ as in
Theorem~\ref{thm:v87-moving-envelope}, define
\begin{equation}\label{eq:v87-exceptional-source}
 R_h^\varepsilon(t)=\int W^\varepsilon(x)
   \one_{\{z(x)<a(h)\ \mathrm{or}\ \ell(x)<\delta(h)\}}\,d\mu_t(x),
 \qquad
 \mathfrak R_{\varepsilon,M}
     =\lim_{h\downarrow0}\mathcal H_M(R_h^\varepsilon).
\end{equation}
This limit exists in $[0,\infty]$: the source sets decrease with $h$.
The two alternatives are a union, not a double-counted sum.
Every boundary type, including grazing forward ends, is retained in
this original-volume restriction. In particular it is not identical
to the earlier source $q_\delta$, which selected nongrazing forward ends.
Cofinality shows that the same charge is obtained by taking the infimum
over all positive pairs of gradient and total-residence thresholds,
always after applying $\mathcal H_M$.

The fixed-positive-residence surface theorem retained in Supplement
Theorem~\ref{supp-v85archive:thm:v85-cutoff-surface} states
\begin{equation}\label{eq:v87-positive-residence-input}
 \mathcal H_M(\overline{\mathcal F}_h^{\ge\delta})
                          \le C h^{1/32}\qquad(h,\delta>0\text{ fixed}).
\end{equation}
Its leading constant is independent of the positive $\delta$. Its
proof sums finitely many physical residence bands before the collision
limsup. We apply it with $\delta=\delta(h)$ for each fixed $h$, not
at $\delta=0$ or at a collision-dependent cutoff.

\begin{theorem}[The original bad-source height is localized]
\label{thm:v87-original-localization}
For every fixed sufficiently small $\varepsilon>0$ and $M$,
\begin{equation}\label{eq:v87-localization-sandwich}
 \mathfrak R_{\varepsilon,M}\le\mathcal H_M(b^\varepsilon)
          \le\mathfrak R_{\varepsilon,M}+C\varepsilon^{1/16}.
\end{equation}
More precisely, for each sufficiently small fixed $h>0$,
\begin{align}
 \mathcal H_M(b^\varepsilon)\le{}&
  \mathcal H_M(R_h^\varepsilon)+C\varepsilon^{1/16}
      +Ch^{1/32}+Kh^{31/1408}\notag\\
 &\hspace{20mm}+C_\Theta(1+\varepsilon^{-1})\sqrt h.
                                         \label{eq:v87-finite-width-source}
\end{align}
All sources use the original section normalization and exact labels.
\end{theorem}
\begin{proof}
Fix the outer width $h$ and its two initial cutoffs. Split the starting
fiber into the exceptional set in \eqref{eq:v87-exceptional-source}
and its complement. For every $0<s<h$, the exceptional contribution
to $E_s^\varepsilon+O_s^\varepsilon$ is at most $R_h^\varepsilon$:
exit and survivor are disjoint cases, and the positive inverse-twist
loss on a survivor is at most its original weight.

On the complementary initial family, negative optical variation pays
surviving loss and grazing first exits as in
Theorem~\ref{thm:v87-moving-envelope}. For a nongrazing first exit,
$\omega(t)\le\omega(b-)+\int_t^b[-\omega']_+$. Its endpoint term,
after dividing by the initial angular normalization, is bounded by a
fixed multiple of the actual outgoing flux on that same component.
The full segment has residence at least $\delta(h)$, so dropping
only the initial-gradient selector on this positive upper term gives
$\mathcal F_s^{\ge\delta(h)}$. The variation is bounded by
$C\mathcal V_h^{a(h),\delta(h)}$. Thus, before any limit,
\[
 \bar E_h^\varepsilon+\bar O_h^\varepsilon
 \le R_h^\varepsilon+C\overline{\mathcal F}_h^{\ge\delta(h)}
                         +C\mathcal V_h^{a(h),\delta(h)}.
\]
The averaging uses the same initial cutoffs throughout $0<s<h$.
The inverse physical Jacobian is paid by the optical variation;
it has not been deleted from the first-exit formula.

Insert this inequality into \eqref{eq:v83-shift-input}. The fixed-width
bad-source estimate bounds $2\mathcal H_M(A_hb^\varepsilon)$ by
$C\varepsilon^{1/16}$, and its unweighted version bounds the guard
term by $C_\Theta(1+\varepsilon^{-1})\sqrt h$.
Equation~\eqref{eq:v87-positive-residence-input} and the finite-count
envelope give \eqref{eq:v87-finite-width-source}. For the latter one
may use one fixed auxiliary band, so $K$ does not change with $h$.
All other auxiliary bands are enlarged only after the collision limit.

Since $0\le R_h^\varepsilon\le b^\varepsilon$, the first inequality
in \eqref{eq:v87-localization-sandwich} is immediate. For the second,
let $h\downarrow0$ in \eqref{eq:v87-finite-width-source} only after
its ordered height has been taken, with $\varepsilon$ fixed.
This argument remains valid for extended nonnegative heights.
\end{proof}

\begin{theorem}[An exact original-source endpoint criterion]
\label{thm:v87-endpoint}
With $\varepsilon(B)=A_0B^{-1/12}$ and the unchanged scalar remainder,
\begin{equation}\label{eq:v87-endpoint-sandwich}
 \mathfrak R_{\varepsilon(B),M}\le\mathcal E_M+C_MB^{-1/192},\qquad
 \mathcal E_M\le\mathfrak R_{\varepsilon(B),M}+C_MB^{-1/192}.
\end{equation}
In particular
\begin{equation}\label{eq:v87-endpoint-equality}
             \mathcal E_M=\lim_{B\to\infty}\mathfrak R_{\varepsilon(B),M}.
\end{equation}
The original full pointwise theorem is equivalent to the right side
being zero for every $M$. This zero-charge assertion is not proved
by the preceding trace and envelope theorems.
\end{theorem}
\begin{proof}
The positive remainder \eqref{eq:v76-remainder} and the triangle
inequality give both
\[
 \mathcal H_M(b^{\varepsilon(B)})\le\mathcal E_M+C_MB^{-1/192},
 \qquad
 \mathcal E_M\le\mathcal H_M(b^{\varepsilon(B)})+C_MB^{-1/192}.
\]
Combine them with \eqref{eq:v87-localization-sandwich} and
$\varepsilon(B)^{1/16}=A_0^{1/16}B^{-1/192}$.
This proves \eqref{eq:v87-endpoint-sandwich} and then
\eqref{eq:v87-endpoint-equality}, also when the value is infinite.
No monotonicity of the original smooth function $\Theta$ in its
argument, or of the source in $B$, is needed for the squeeze.
For every fixed $B$, all physical widths used above are fixed before
the collision limit. The width tends to zero afterwards; $B$ is last.
No arithmetic class is normalized in this argument.
\end{proof}

The prior sufficient pairs $(\mathfrak C_M,\mathfrak B_M)$ and
$(\mathfrak C_M,\mathfrak Q_M)$ remain valid in the preserved article.
The new criterion is an equivalence for a restriction of the original
bad source, not an assertion that either earlier, larger optical charge
has vanished. It isolates collapsing initial gradient and collapsing
total residence while estimating their complementary moving family
uniformly at finite count. Source-specific control of the remaining
restriction is still required. Scalar closure must precede every
same-roof path ratio and every unconditional likelihood consequence.
